\subsection{问题四分析}

第四问的目标不是用新增实验重复证明前三问，而是在5次预算内优先减少最影响条件选择的
不确定性。前三问暴露的数据不足可以归为三类：关键温度区间观测稀疏、最高观测缺少
独立重复、不同组合缺少新增同温条件下的直接比较。因此，实验设计既要补充局部温度
信息，也要兼顾可复现性和竞争组合比较，不能只按某个拟合模型的预测收率排序。

要刻画各组合在整个温度范围内的弯曲变化，即建立完整响应面，需要更多实验；本题的
催化剂组合又是不可随意拆分的离散类别。5次预算不足以同时确定全温区的曲线形态、
组合差异和随机误差。本文因而采用证据缺口驱动的局部增补：
先由前三问已经得到的结果明确每项未决问题，再为每次实验规定明确且可核查的职责。
具体缺口、实验条件和观测后的更新规则均在第四问模型建立与求解中给出
\cite{nistreplication,nistrsmcompare}。
